% RESEARCH_DRAFT_SHA256: 45b99f210a5b8e39051d0ba0e75a3c083db89890d9bc67be1ccd4a1f1d69d7e3
% Layout and source-reference wrappers expanded for the web renderer.
\documentclass[11pt]{article}
\usepackage[margin=25mm]{geometry}
\usepackage{amsthm}
\usepackage{fontspec,unicode-math}
\setmainfont{Latin Modern Roman}
\setmathfont{Latin Modern Math}
\usepackage{amsmath,amssymb,mathtools,xurl,hyperref}
\hypersetup{colorlinks=true,urlcolor=blue,linkcolor=blue}
\setcounter{secnumdepth}{0}
\setlength{\parindent}{0pt}
\setlength{\parskip}{5pt}
\setlength{\emergencystretch}{3em}
\newcommand{\sref}[2]{\hyperlink{source:#1:#2}{[S#2]}}
\newcommand{\cataloguescope}[1]{\paragraph{Source record.}#1}
\newcommand{\floor}[1]{\left\lfloor#1\right\rfloor}
\newtheorem{proposition}{Proposition}
\newtheorem{lemma}[proposition]{Lemma}
\begin{document}
\paragraph{Research manuscript.} OpenAI \texttt{gpt-6-astra}, reasoning \texttt{ultra}; 11 October 2026.
\paragraph{Preserved source body.} The complete catalogue mathematical text and sources follow with layout and source links expanded, including source qualifications and its historical preparation notes. New research begins at the section headed Research attempt.

\section{Machine-Learning Algorithms Results Report}\label{statement}

\subsection{Definitions and mathematical statement}

\paragraph{Shared notation.}Standard number systems are $\mathbb N,\mathbb Z,\mathbb Q,\mathbb R,\mathbb C$. Specialized notation and index ranges follow the cited source and the definitions below.


A learning problem specifies a data distribution, a hypothesis class and a loss. Statistical consistency means that the estimator approaches its population target as the sample size increases; statistical and computational requirements are distinct.
\paragraph{Record scope and qualifications.}
The catalogue title is retained as metadata. Named certificates, contracts, bounds or frameworks in that title are not asserted as results of the original source extract.
The supplied registry flags this extraction as unrecoverable. It is displayed as a literal source fragment; the definitions provide context rather than a replacement formal statement.

Source-record limitation: Source flags the statement as unrecoverable: The official contribution asks rhetorically whether provable algorithms are possible and immediately answers yes by presenting three results; the extraction is expository framing, not a residual open problem. No separate clean statement is promoted; the original record remains publicly visible and retained for provenance.
The source text below is rendered literally where necessary; damaged formulas and composite extracts are not silently repaired into a new mathematical claim.
The supplied catalogue marks this record as unrecoverable. Its extract and terminology are documented, but a complete original problem has not been reconstructed.
\paragraph{Statement recorded in the supplied source.}
Machine-learning tasks are often computationally intractable, while practical algorithms lack provable guarantees on running time or solution quality. Can one develop algorithms for such tasks with provable running-time and solution-quality guarantees? \href{https://doi.org/10.4171/owr/2012/54}{[S2]}
\paragraph{Source status.}
The supplied catalogue labels this source record open; that metadata was checked against the supplied snapshot on 2026-10-11. The label is preserved for provenance and is not an independent verification of present open status or mathematical correctness. The catalogue literature assessment is dated 2026-08-21.
\subsection{Short English statement}
The supplied fragment concerns machine-Learning Algorithms Results Report. Its complete mathematical target is not recovered.
\subsection{Sources}
\begin{itemize}
\item[S1] Ulam AI and the UnsolvedMath Contributors, UnsolvedMath: A Curated Collection of Open Mathematics Problems, catalogue entry. Catalogue curation and metadata: CC BY 4.0. The source extract is retained; definitions, English summaries and qualifications are editorial additions.
\url{https://huggingface.co/datasets/ulamai/UnsolvedMath}.
\item[S2] Complexity Theory (2013), Oberwolfach Reports, DOI 10.4171/owr/2012/54, pp. 3267--3304
\url{https://doi.org/10.4171/owr/2012/54}.
\item Source-specific Oberwolfach Reports licences, including CC BY-SA 4.0 where applicable, are retained; see THIRD\_PARTY\_NOTICES.md.
\end{itemize}
\paragraph{Preparation and review.}
The original source and UnsolvedMath contributors retain attribution for the source extract. Codex (GPT-6; exact serving revision not exposed) assisted with definitions, English summaries, scope notes and formatting. No new human mathematical authorship or independently verified solution is asserted. Catalogue identity checked against the complete supplied registry and current repository. Distinct upstream records, including repeated questions and broad or damaged extracts, are retained for complete coverage of the source CS category. A repeated record is not a new independent conjecture. Complete mathematical-equivalence and current-literature review has not been performed. Preparation and final review follow contribution prompts C and D, with source gaps disclosed. Structural and rendering checks do not establish a new theorem or open status.
\label{end:statement}

\section{Research attempt}
\subsection{Literature, source context and scope}
The full original catalogue body and source qualifications are preserved above. The relevant primary workshop passages were read in the cached publisher report. Historical claims and subsequent attributed results are distinguished from the mathematical work below. A restricted example, standard lemma or known theorem is not counted as new progress on the entire original target. Missing assumptions and repeated or rhetorical extracts are not silently turned into different solved problems.

\subsection{A small class can be learned by exhaustive search}
The original talk already gives positive examples in machine learning and asks a broad research-program question. We test a restricted explanation for when a theoretically justified learner is efficient. Let $\mathcal H$ be an explicitly enumerable finite family of $M\ge1$ binary functions. Assume the target $h_*$ belongs to this family, training inputs are independent with common distribution $P$, and labels are exactly $h_*(x)$. Choose any hypothesis having zero training error.
\begin{proposition}
For $0<\alpha,\delta<1$, if $n\ge(\log M+\log(1/\delta))/\alpha$, every training-consistent hypothesis has true error at most $\alpha$ with probability at least $1-\delta$. Exhaustive search costs at most $nM$ hypothesis evaluations.
\end{proposition}
\begin{proof}
Fix $h\in\mathcal H$ of true error $p>\alpha$. Independence gives probability $(1-p)^n\le e^{-\alpha n}$ that it is consistent with all $n$ labels. Taking the union bound over at most $M$ hypotheses bounds the probability of any bad consistent one by $Me^{-\alpha n}\le\delta$. A consistent hypothesis always exists because $h_*\in\mathcal H$. Testing every hypothesis on every example proves the cost bound.
\end{proof}
For a worked case take a uniform two-point input space, target identically zero, and all four binary functions as $\mathcal H$. For $n\ge1$, a hypothesis with positive error remains consistent exactly when at least one input point was never sampled. These two missing-point events are disjoint and have total probability $2^{1-n}$. This can be checked by listing all $2^n$ input sequences.
The proposed shortcut from a logarithmic sample dependence on $M$ to a logarithmic running time fails: the proof uses an exhaustive search costing $M$, and supplies no efficient optimization oracle. If $M=2^d$, that cost is exponential in $d$. No general lower bound is claimed; particular classes can have much better algorithms.
\subsection{Verification and remaining gaps}
All input sequences for $1\le n\le10$ were enumerated, reproducing $2^{1-n}$ exactly. The realization, independence, and explicit-enumeration hypotheses were checked separately. This is a standard finite-class learning bound, not a novel answer to the broad workshop agenda. Noise, misspecification, implicit large classes and empirical performance remain outside this attempt.

\section{References and source attribution}
Primary workshop report: \url{https://doi.org/10.4171/owr/2012/54}. Relevant passages read: 2012/54 pp. 3298--3299, Arora’s question and existing positive examples from topic models, nonnegative factorization and independent component analysis.

All retrievals were checked on 11 October 2026. Original question attribution remains in the preserved Sources section. External theorems are attributed and not claimed as new discoveries or fully verified merely by citation.


\paragraph{Integration review.} Codex (GPT-6; exact serving revision not exposed) reviewed the written arguments and expanded layout and source-reference wrappers. Repository metadata and the production web renderer were checked using Python 3.13 and JavaScript with MathJax 3.2.2. This model-assisted review is not independent human proof verification; no PDF compilation or formal proof checking is asserted.

\section{Completion Estimate}
\noindent\textbf{COMPLETION: 0\%}
No substantiated new progress toward the entire original target is claimed. The worked results and failed route are delimited above; a known result, illustrative model, or auxiliary obstruction is not counted as original-target completion.

\section{Final status and attribution}
\textbf{Attempt status: unresolved.} This describes this research attempt, not a new ruling on the source catalogue's status or on external literature claims.
The mathematics and this self-audit were generated by OpenAI \texttt{gpt-6-astra}, reasoning \texttt{ultra}, on 11 October 2026. Codex (GPT-6; exact serving revision not exposed) orchestrated the batch, reviewed the written mathematical arguments, and contributed formatting/repository review. No human mathematical authorship, independent human proof checking, or formal verification is claimed. The historical preparation paragraphs in the reproduced catalogue body concern that body only.

Contribution prompt A was applied to the complete definition and supplied source record, followed by the adversarial checks of prompt B. Local reproduction files are \texttt{batch019/checks.py}, \texttt{computations.json}, and the per-ID audit JSON. They were not public links at preparation time. Numerical checks are finite and do not prove unrestricted results. No LaTeX engine was found locally; compilation was not performed by this research agent.

\end{document}
